% Save as: syllabus.tex
% ============================================================================
% Course Syllabus — Microeconomic Theory for Experimentalists & Applied
% Microeconomists: Foundations, Behavior & AI-Augmented Experiments
% 14-week semester version. Compile with pdflatex.
% ============================================================================
\documentclass[11pt]{article}
\usepackage[margin=1in]{geometry}
\usepackage{amsmath, amssymb}
\usepackage{booktabs}
\usepackage{enumitem}
\usepackage{longtable}
\usepackage[hidelinks]{hyperref}

\title{Microeconomic Theory for Experimentalists \& Applied
Microeconomists\\
\large Foundations, Behavior \& AI-Augmented Experiments}
\author{\textbf{[Instructor]} \quad · \quad \textbf{[email]} \quad ·
\quad Office hours: \textbf{[times/location]}}
\date{\textbf{[Semester Year]} · \textbf{[meeting times]} ·
\textbf{[room]}}

\begin{document}
\maketitle

\section*{Course description}
This course trains master's students to \emph{design, critique, and
pilot real experiments}. It covers the microeconomic theory an
experimentalist actually uses --- choice under uncertainty, consumer
and producer theory, information economics, behavioral economics, and
causal inference --- taught intuition-first: through graphs, worked
numerical examples, and in-class demonstrations in which you generate
the anomalies with your own choices before the theory names them.

The course's distinguishing feature is serious integration of LLM
generative agents (``Homo silicus''-style simulation) as both a
pedagogical tool and a research method: every one to two weeks you run
a simulation lab in which populations of AI agents play the games and
face the decisions the theory describes, and you compare their behavior
to the human experimental literature --- and, often, to this class's
own in-class choices. The comparison is treated as a first-class
intellectual question, never a gimmick: by the end you will know
precisely what silicon pilots can establish (instrument debugging,
pipeline validation, power analysis, framing detection, audits of
deployed AI systems) and what they cannot (effect-size forecasts,
welfare claims), and you will have the habits --- preregistration,
stability audits, honest methods paragraphs --- that separate the two.

\textbf{No prior programming experience is assumed.} All lab code is
provided complete and runs top-to-bottom in Google Colab; your
modifications are confined to clearly marked \texttt{CHANGE HERE}
blocks, and every notebook ships with a \texttt{DRY\_RUN} mode that
executes end-to-end without an API key or cost.

\section*{Learning objectives: the three exit competencies}
Graduates of this course can:
\begin{enumerate}[nosep]
  \item \textbf{Read modern applied and experimental papers
        critically} --- identify what was tested, whether the design
        supports the claims, and what would strengthen it.
  \item \textbf{Design theory-informed experiments end to end} ---
        subjects, treatments, incentive-compatible elicitation,
        identification, simulation-based power, and ethics as design
        substance. \emph{This competency carries the most weight.}
  \item \textbf{Use LLM generative agents to pilot, stress-test, and
        scale designs} before spending money on human subjects ---
        with explicit, disciplined reasoning about where silicon
        subjects diverge from humans.
\end{enumerate}

\section*{Format}
\begin{itemize}[nosep]
  \item \textbf{Lectures} --- theory plus the story of its application
        (classic experiments and recent papers), with live in-class
        choice demonstrations wherever the material allows. Expect to
        make paid decisions on your phone most weeks.
  \item \textbf{AI tutor sessions} (weekly, required) --- structured
        conversations with a course-configured tutor AI that quizzes
        you, debugs your experiment designs, and adapts to your
        research interests. You submit the session log with a short
        reflection.
  \item \textbf{LLM simulation labs} (every 1--2 weeks) --- run and
        modify generative-agent experiments; compare agent behavior to
        human benchmarks; write up findings in the course's
        six-section format.
  \item \textbf{Experiment sessions} --- in weeks 11--13 the second
        meeting is a hands-on session: a calibration workshop, a
        common-pool resource experiment, and a pricing tournament,
        each played for real (small) stakes.
  \item \textbf{Capstone} --- design and simulate an original
        experiment; present a proposal ready for an IRB, an employer,
        or a PhD application. See \texttt{capstone/} in the course
        repository.
\end{itemize}

\section*{Assessment}
\begin{center}
\begin{tabular}{@{}lc@{}}
\toprule
Component & Weight \\
\midrule
Weekly AI tutor logs + reflections & 10\% \\
Problem sets (10; lowest dropped) & 25\% \\
Simulation-lab write-ups & 15\% \\
Midterm: replicate + extend a published experiment & 20\% \\
Capstone (milestones 10\% · paper 60\% · pilot 15\% ·
presentation 15\% of this component) & 30\% \\
\bottomrule
\end{tabular}
\end{center}

\noindent\emph{Problem sets} mix guided calculations, ``design an
experiment to test this prediction'' tasks (the heaviest-weighted
part of every set), and silicon-pilot interpretation. \emph{Lab
write-ups} use the required six-section structure (question/benchmark
$\to$ agent design $\to$ treatments $\to$ results vs.\ humans $\to$
divergence diagnosis $\to$ design implication). The \emph{midterm}
(due week 9) asks you to replicate and extend one published experiment
using either human data from the paper or an LLM simulation, in a
short paper with the same six-section discipline.

\section*{Weekly schedule}
Core readings are read deeply (2--3 items per week); optional readings
are for going further and for capstone mining. Full reading details,
lab instructions, and deliverable checklists are in each module's
student handout in the course repository.

\begin{longtable}{@{}p{0.7cm}p{4.3cm}p{5.6cm}p{3.6cm}@{}}
\toprule
Wk & Module & Core readings & Due \\
\midrule
\endhead
1--2 & \textbf{1.} Rational Choice, Preferences \& Uncertainty &
Kahneman \& Tversky (1979); Holt \& Laury (2002) &
Lab 1 write-up; PS1 \\
3 & \textbf{2.} Consumer Theory, Demand \& Welfare &
Kahneman, Knetsch \& Thaler (1990); Andreoni \& Miller (2002) &
Lab 2 write-up; PS2 \\
4 & \textbf{3.} Producer Theory, Costs \& Technology &
Duflo, Kremer \& Robinson (2011); Bloom et al.\ (2013) &
Lab 3 write-up; PS3 \\
5--6 & \textbf{4.} Asymmetric Information, Moral Hazard \& Adverse
Selection & Fehr, Kirchsteiger \& Riedl (1993); Lazear (2000) &
Lab 4 write-up; PS4 \\
7--8 & \textbf{5.} Behavioral Economics: Deviations, Biases \& Nudges &
Kahneman \& Tversky (1979) revisited; Ashraf, Karlan \& Yin (2006);
Madrian \& Shea (2001) &
Lab 5 write-up; PS5; \emph{capstone question memo (wk 8)} \\
9--10 & \textbf{6.} Causal Inference, RCTs \& Structural Integration &
Bertrand \& Mullainathan (2004); Duflo, Glennerster \& Kremer (2007) &
\emph{Midterm (wk 9)}; Lab 6 write-up; PS6; \emph{capstone theory +
hypotheses (wk 10)} \\
11 & \textbf{7.} Frontier: Development, Labor \& Education ·
\emph{calibration workshop} & Cunha \& Heckman (2007); Miguel \&
Kremer (2004) & Lab 7 write-up; PS7 \\
12 & \textbf{8.} Frontier: Environmental, Resource \& Conservation ·
\emph{CPR experiment session} & Ostrom, Walker \& Gardner (1992);
Jayachandran et al.\ (2017) & Lab 8 write-up; PS8; \emph{capstone
design draft} \\
13 & \textbf{9.} Frontier: Digital \& Platform Markets ·
\emph{pricing tournament} & Calvano et al.\ (2020); Blake, Nosko \&
Tadelis (2015) & Lab 9 write-up; PS9 \\
14 & \textbf{10.} AI-Econ Nexus \& Capstone ·
\emph{capstone workshop} & Horton (2023); Argyle et al.\ (2023) &
Lab 10 report; PS10; \emph{capstone simulation pilot +
preregistration} \\
15 & Capstone & --- & \emph{Comparison memo} \\
F & Capstone presentations & --- & \emph{Final paper +
presentation} \\
\bottomrule
\end{longtable}

\noindent Optional/going-further readings for every module (classics
such as Allais 1953, Ellsberg 1961, Holmstr\"om 1979, Rochet \& Tirole
2003; surveys such as DellaVigna 2009 and Foster \& Rosenzweig 2010;
and one recent LLM-agents paper per module, announced in class) are
listed in the module handouts. Textbook support: Nicholson \& Snyder,
\emph{Microeconomic Theory} (theory backbone, chapters assigned per
module); Kagel \& Roth, \emph{Handbook of Experimental Economics}
(selected chapters); Thaler \& Sunstein, \emph{Nudge} (behavioral
modules, mechanism taxonomy).

\section*{Software and materials}
\begin{itemize}[nosep]
  \item \textbf{Google Colab} (free tier sufficient) --- all labs run
        here; nothing is installed locally.
  \item \textbf{Anthropic API key} --- required for live lab runs; a
        modest budget (typically under \$5 per lab) suffices, and
        every notebook's \texttt{DRY\_RUN} mode works without any key.
        \textbf{[Instructor: state here whether course API credits are
        provided.]}
  \item \textbf{Course repository} ---
        \url{https://github.com/moncap/micro-course}: labs (with
        Open-in-Colab links), problem sets, handouts, and capstone
        materials.
  \item \textbf{LaTeX} (Overleaf recommended) --- problem-set and
        capstone write-ups may be submitted in LaTeX or Word; LaTeX
        templates are provided and its use is encouraged.
  \item \textbf{Expected Parrot EDSL} (optional) --- for
        survey-style capstone simulations; plain-Python scaffolds
        cover everything required.
  \item In-class experiments run on \textbf{oTree}/Qualtrics
        instances operated by the instructor; you only need a phone or
        laptop with a browser.
\end{itemize}

\section*{AI use policy}
This course \emph{requires} AI use --- and demands discipline about it.
\begin{itemize}[nosep]
  \item \textbf{Required:} the weekly tutor sessions, the simulation
        labs, and the capstone pilot are AI-based coursework. Problem
        sets flag AI-permitted parts with $\langle$AI$\rangle$.
  \item \textbf{Disclosure:} any substantive AI assistance ---
        required or voluntary --- must be disclosed, and transcripts
        attached where the assignment says so. Keep your logs; they
        are part of your grade and your protection.
  \item \textbf{Boundaries:} AI may draft code within the labs'
        \texttt{CHANGE HERE} conventions and may serve as a study
        partner anywhere. The analysis, interpretation, and prose of
        your write-ups must be your own; verbatim AI-written analysis
        submitted as yours is an academic-integrity violation.
  \item \textbf{Provenance:} simulated data must never be presented as
        human data, in this course or anywhere else. This is an
        integrity constraint, not a formality.
  \item \textbf{Honesty standard:} claims about what a simulation
        shows are graded against the course's can/cannot framework;
        over-claiming costs more than admitting a limitation.
\end{itemize}

\section*{Course policies}
\begin{itemize}[nosep]
  \item \textbf{Collaboration:} discussing ideas on problem sets is
        encouraged; write-ups are individual. Lab write-ups and the
        capstone follow their own stated rules (capstone pairs
        allowed with a contribution statement).
  \item \textbf{Late work:} lab write-ups and problem sets lose 10\%
        per day, capped at 50\%; the lowest problem set is dropped.
        Capstone milestones are gating --- late milestones delay
        feedback, not points, but the final deliverables are firm
        dates.
  \item \textbf{In-class experiments} pay real (small) stakes and are
        part of the course's data; participation is expected, and any
        student may opt out of having their (always anonymized)
        choices displayed by telling the instructor.
  \item \textbf{Accessibility:} students needing accommodations should
        contact \textbf{[disability services office]} and the
        instructor in week 1.
  \item \textbf{Academic integrity:} governed by \textbf{[institutional
        policy]}; the AI-use policy above defines the line for
        AI-related work.
\end{itemize}

\section*{One paragraph of advice}
The students who thrive in this course are not the strongest coders or
the most mathematically fluent --- the code is given and the math is
scaffolded. They are the ones who take the design questions seriously:
who ask, of every model, ``what experiment does this let me build?''
and, of every simulation, ``what exactly is this evidence of?'' Bring
your own research interests --- the tutor sessions, the problem sets'
design parts, and the capstone are all built to run on them.

\end{document}
